% from ega2-errata-addenda-fr.tex
\oldpage[II]{220}
\markboth{A. GROTHENDIECK --- Chap. II}{A. GROTHENDIECK --- Chap. II}
annelés. Pour tout $A$-module $M$, on a (avec les notations de (1.8.4)) un
isomorphisme canonique fonctoriel de $\sh{O}_X$-Modules
\begin{equation}
\label{ega2:addendum:I.1.8.6.1-fr}
  u^*(\widetilde M)\isoto u^*(i^*(M')).
\tag{1.8.6.1}
\end{equation}
\end{corollary}

\begin{corollary}[1.8.7]
\label{ega2:addendum:I.1.8.7-fr}
Sous les hypothèses de (1.8.6), on a, pour tout $A$-module $M$ et tout
$\sh{O}_X$-Module $\sh{F}$, un isomorphisme canonique fonctoriel en
$M$ et $\sh{F}$
\begin{equation}
\label{ega2:addendum:I.1.8.7.1-fr}
  \Hom_{\sh{O}_S}(\widetilde M,u_*(\sh{F}))
  \isoto\Hom_A(M,\Gamma(X,\sh{F})).
\tag{1.8.7.1}
\end{equation}
\end{corollary}

On a en effet, en vertu de (0, 4.4.3) et du lemme (1.8.5), un
isomorphisme canonique de bifoncteurs
\[
  \Hom_{\sh{O}_S}(\widetilde M,u_*(\sh{F}))
  \isoto\Hom_{A'}(M',i_*(u_*(\sh{F}))),
\]
et il est clair que le second membre n'est autre que
$\Hom_A(M,\Gamma(X,\sh{F}))$. On notera que l'homomorphisme canonique
(1.8.7.1) fait correspondre à tout $\sh{O}_S$-homomorphisme
$h:\widetilde M\to u_*(\sh{F})$ (autrement dit, à tout $u$-morphisme
$\widetilde M\to\sh{F}$) le $A$-homomorphisme
\[
  \Gamma(h):M\to\Gamma(S,u_*(\sh{F}))=\Gamma(X,\sh{F}).
\]

\begin{env}[1.8.8]
\label{ega2:addendum:I.1.8.8-fr}
Avec les notations de (1.8.4), il est clair (0, 4.1.1) que tout
morphisme d'espaces annelés $(\psi,\theta):X\to S'$ équivaut à la
donnée d'un homomorphisme d'anneaux
$A\to\Gamma(X,\sh{O}_X)$. On peut donc encore interpréter (1.8.1)
comme définissant une bijection canonique
\[
  \Hom(X,S)\isoto\Hom(X,S')
\]
(où bien entendu il s'agit au second membre de morphismes d'espaces
annelés, puisqu'en général $A$ n'est pas un anneau local). Plus
généralement, si $X,Y$ sont deux espaces annelés en anneaux locaux et
si $(Y',A')$ est l'espace annelé dont l'espace sous-jacent est réduit à
un point et dont le faisceau d'anneaux $A'$ est le faisceau simple
défini par l'anneau $\Gamma(Y,\sh{O}_Y)$, on peut interpréter
(1.8.2.1) comme une application
\begin{equation}
\label{ega2:addendum:I.1.8.8.1-fr}
  \rho:\Hom_{\mathrm{an}}(X,Y)\to\Hom(X,Y').
\tag{1.8.8.1}
\end{equation}
Le résultat de (1.8.3) s'interprète donc en disant que les schémas
affines sont caractérisés parmi les espaces annelés en anneaux locaux
comme ceux pour lesquels la restriction de $\rho$ à $\Hom(X,Y)$ :
\begin{equation}
\label{ega2:addendum:I.1.8.8.2-fr}
  \rho':\Hom(X,Y)\to\Hom(X,Y')
\tag{1.8.8.2}
\end{equation}
est bijective pour tout espace annelé en anneaux locaux $X$. Dans un
chapitre ultérieur, nous généraliserons cette définition, ce qui
permettra d'associer à un espace annelé quelconque $Z$ (et non plus
seulement à un espace annelé dont l'espace sous-jacent est réduit à un
point) un espace annelé en anneaux locaux que nous appellerons encore
$\Spec(Z)$; cela sera le point de départ d'une théorie «~relative~» des
préschémas au-dessus d'espaces annelés quelconques, étendant les
résultats du chap.~I.
\end{env}

\begin{env}[1.8.9]
\label{ega2:addendum:I.1.8.9-fr}
On peut considérer les couples $(X,\sh{F})$ formés d'un espace annelé
en anneaux locaux $X$ et d'un $\sh{O}_X$-Module $\sh{F}$ comme formant
une catégorie, un morphisme $(X,\sh{F})\to(Y,\sh{G})$ de cette catégorie
étant un couple $(u,h)$ formé d'un morphisme d'espaces
